\section{Numerical Model}
\label{sec:numerical-model}

The simulations use an effective driven two-level system coupled to Markovian
relaxation and dephasing channels.  In the frame rotating at the drive
frequency, and within the rotating-wave approximation, the Hamiltonian is
implemented as
\begin{equation}
    \frac{H(t)}{\hbar}
    =
    \Delta a^\dagger a
    + \frac{\Omega_0}{2} f(t)
      \left(a+a^\dagger\right),
    \label{eq:numerical-hamiltonian}
\end{equation}
where $\Delta=\omega_q-\omega_d$ is the qubit--drive detuning, $\Omega_0$ is
the peak Rabi frequency, and $a=|0\rangle\langle1|$ is the lowering operator.
Up to an irrelevant identity term and the chosen sign convention for
$\Delta$, Eq.~\eqref{eq:numerical-hamiltonian} is equivalent to
$H/\hbar=(\Delta/2)\sigma_z+(\Omega_0 f/2)\sigma_x$.

For a root-Lorentzian pulse of total duration $L$, the dimensionless envelope
is
\begin{equation}
    f_{\mathrm{RL}}(t)
    =
    \left[1+\left(\frac{t}{\tau}\right)^2\right]^{-1/2},
    \qquad -\frac{L}{2}\le t\le\frac{L}{2},
\end{equation}
with $\tau$ chosen such that
$f_{\mathrm{RL}}(\pm L/2)=c$, where $c$ is the relative cutoff.  The echo
variant is obtained by changing the drive phase by $\pi$ at the pulse center,
\begin{equation}
    f_{\mathrm{ERL}}(t)
    = \operatorname{sgn}(-t) f_{\mathrm{RL}}(t).
    \label{eq:numerical-echo-envelope}
\end{equation}
The overall sign in Eq.~\eqref{eq:numerical-echo-envelope} is a phase
convention and does not affect the measured populations.

The density matrix evolves according to the Lindblad master equation
\begin{equation}
    \dot{\rho}
    =
    -\frac{i}{\hbar}[H(t),\rho]
    + \mathcal{D}\!\left[\sqrt{\frac{1}{T_1}}\,a\right]\rho
    + \mathcal{D}\!\left[\sqrt{\frac{2}{T_\phi}}\,
      a^\dagger a\right]\rho,
    \label{eq:numerical-master-equation}
\end{equation}
where
$\mathcal{D}[C]\rho=C\rho C^\dagger-\frac{1}{2}
\{C^\dagger C,\rho\}$ and
\begin{equation}
    \frac{1}{T_2}=\frac{1}{2T_1}+\frac{1}{T_\phi}.
\end{equation}
Each calculation starts in $|0\rangle$, integrates
Eq.~\eqref{eq:numerical-master-equation} over the complete pulse, and records
the final Bloch-vector components.  Spectroscopy maps are generated by
repeating this evolution over the experimental detuning and drive-amplitude
grids, using the pulse duration, cutoff, $T_1$, and $T_2$ associated with the
corresponding data set.

The numerical Hilbert space is restricted to two levels.  Consequently, this
model captures coherent drive dynamics, relaxation, and pure dephasing, but it
does not describe leakage, AC-Stark shifts produced by higher transmon levels,
or multiphoton transitions.  Comparisons in the extended-amplitude regime are
therefore interpreted as tests of where the effective two-level description
breaks down rather than as a simulation of those multilevel effects.
